\section{数值微分为何不能无限缩小步长}
\label{sec:08-Numerical-Analysis/04-Numerical-Differentiation-and-Integration/01}

学过微积分的人看到导数定义

\[
f'(x)=\lim_{h\to0}\frac{f(x+h)-f(x)}{h},
\]

脑子里会冒出一个很自然的程序思路：取一个极小的 \(h\)，算一下差分，除以 \(h\)，完事。如果 \(10^{-6}\) 不够精确，换 \(10^{-12}\)。再不行就换 \(10^{-16}\)——双精度的极限。极限说的是 \(h\to0\)，那程序就应该尽量逼近零，不是吗？

不是。我们用具体的数字来试。取 \(f(x)=\sin x\)，求 \(x=1\) 处的导数。真值是 \(\cos 1\approx 0.54030230586\)。前向差分的结果：

\[
\begin{array}{c|c|c}
h & \text{近似值} & \text{误差} \\
\hline
10^{-1} & 0.49736 & 4.3\times10^{-2} \\
10^{-2} & 0.53608 & 4.2\times10^{-3} \\
10^{-4} & 0.54030 & 2.6\times10^{-6} \\
10^{-6} & 0.540302 & 7.3\times10^{-9} \\
10^{-8} & 0.5403023 & \mathbf{1.4\times10^{-9}} \\
10^{-10} & 0.54030 & 1.8\times10^{-6} \\
10^{-12} & 0.539 & 1.0\times10^{-3} \\
10^{-14} & 0.55 & 1.0\times10^{-2} \\
10^{-16} & 0 & 5.4\times10^{-1} \\
\end{array}
\]

误差先下降，到 \(h\approx 10^{-8}\) 附近触底，之后不仅回升，还加速恶化，最后直接给出零。这不是计算错误——如果你在双精度下计算 \(\sin(1+10^{-16})\)，\(1+10^{-16}\) 距离下一个可表示的浮点数还差约 50 倍，被舍入成了 \(1\)。\(\sin(1+10^{-16})=\sin(1)\)，分子为零，除以任何数都是零。极限定义假设每一步算术都精确，程序没有这个前提。

这个 U 形误差曲线是本节的核心研究对象。两股力量在反方向拉扯：步长大时截断误差主宰，步长小时舍入误差爆发。最优步长不在 \(h\to0\) 的极限处，而在两股力量平衡的谷底。

\subsection{Taylor 展开给出差分公式的截断误差}

你要逼近的是 \(f'(x)\)。把 \(f(x+h)\) 在 \(x\) 处 Taylor 展开：

\[
f(x+h)=f(x)+hf'(x)+\frac{h^2}{2}f''(\xi),
\]

其中 \(\xi\) 在 \(x\) 和 \(x+h\) 之间。移项并除以 \(h\)：

\[
D_h^+f(x)=\frac{f(x+h)-f(x)}{h}
=f'(x)+O(h).
\]

这是前向差分，一阶精度，因为扔掉的项是 \(h\) 的一阶量。每把 \(h\) 减半，截断误差大约减半。

中心差分用到了 \(x\) 两侧的信息。同时展开 \(f(x+h)\) 和 \(f(x-h)\)：

\[
\begin{aligned}
f(x+h)&=f(x)+hf'(x)+\frac{h^2}{2}f''(x)+\frac{h^3}{6}f'''(\xi_+),\\
f(x-h)&=f(x)-hf'(x)+\frac{h^2}{2}f''(x)-\frac{h^3}{6}f'''(\xi_-).
\end{aligned}
\]

两式相减，\(f(x)\) 和 \(f''(x)\) 项因对称性精确抵消（设 \(f\in C^3\)，两个余项的中值点可由介值定理合并为一个 \(\xi\)）：

\[
D_h^0f(x)
=\frac{f(x+h)-f(x-h)}{2h}
=f'(x)+\frac{h^2}{6}f'''(\xi).
\]

中心差分的截断误差从一阶跳到二阶，用了两次函数值。多出来的一次函数值买的不是"多一个信息点"，而是买了对称性。对称性让偶次幂项在相减中自动消掉，留下来的第一个未抵消项就是 \(h^2\) 阶。这是许多数值方法反复出现的主题——通过对称放置样本点，免费获取一个额外精度阶。

在区间边界无法访问 \(x-h\) 时必须用单边公式；若使用超过两个节点，可以通过插值多项式求导构造更高阶的差分模板（stencil）。更高阶的模板减少光滑函数的截断误差，但使用更宽的邻域和更大的正负交替权重——这意味着对噪声和边界效应更敏感。阶数不是免费的午餐。

\subsection{总误差的 U 形：两股力量的交汇}

截断误差是 \(h\) 减小时变好的一方。但浮点计算中还有另一方。每个函数值含绝对舍入误差约 \(C_ru\)，其中 \(u\) 是机器舍入单位——双精度下约 \(1.1\times10^{-16}\)，\(C_r\) 约等于函数值本身的量级。计算差分 \(f(x+h)-f(x)\) 时，两个接近数的相减让有效位数大量丢失。再除以 \(2h\)，残存的舍入误差被 \(1/h\) 放大。

以中心差分为例，总误差的粗略模型是

\[
E(h)\approx C_th^2+\frac{C_ru}{h}.
\]

第一项 \(C_th^2\) 随着 \(h\) 减小而减小——这就是截断误差在改善。第二项 \(C_ru/h\) 随着 \(h\) 减小而增大——舍入相消的代价被小分母放大。两者在相反的斜率上运动，交叉处就是总误差的最小值。

平衡两项：令 \(C_th^2=C_ru/h\)，解得 \(h^3=C_ru/C_t\)，给出主导量级

\[
h\asymp u^{1/3},
\]

其中 \(\asymp\) 表示量级相同，只保留了主导幂次。对于双精度，\(u^{1/3}\approx 4.8\times10^{-6}\)，这解释了为何中心差分的最优步长大约在 \(10^{-5}\) 到 \(10^{-6}\) 量级——不是 \(10^{-16}\)。

前向差分的截断误差是 \(O(h)\) 而非 \(O(h^2)\)，误差模型变为

\[
E(h)\approx C_th+\frac{C_ru}{h}.
\]

平衡得到 \(h\asymp u^{1/2}\)。双精度下约 \(10^{-8}\)——比中心差分的最优步长更小，但绝对误差却更大（因为截断误差随 \(h\) 只线性下降而非平方下降）。中心差分既更精确，又可以承受更大的实用步长——后者在函数的二阶导数取之不易时反而是优势。

这些指数不是无条件默认步长。如果你的函数值来自仿真容差或物理测量，那么有效的"噪声单位"不是机器精度 \(u\) 而是测量噪声的幅度 \(\eta\)——把 \(u\) 换成 \(\eta\)，最优 \(h\) 会显著右移。越粗糙的数据，越不能信任小步长。

\subsection{微分天然放大高频}

设你的观测包含快速振荡的小噪声：

\[
y(x)=f(x)+\varepsilon\sin(\omega x).
\]

噪声的幅度只有 \(\varepsilon\)，看似无害。但噪声的导数是 \(\varepsilon\omega\cos(\omega x)\)——幅度变成了 \(\varepsilon\omega\)。频率 \(\omega\) 越高，微分放大越强。一个幅度 \(10^{-3}\)、频率 \(1000\) 的噪声，在数值导数里会变成幅度 1 的振荡——把真实信号全部淹没。

对含噪数据硬套高阶差分公式，你可能得到一个公式上非常光滑、结果却非常离谱的"导数"。高阶差分假设函数在模板宽度内被低次多项式良好逼近；噪声打破了这一前提。

处理含噪数据的更可靠方向是：把"估计导数"视为一个带正则化的局部建模问题。先用局部多项式拟合、平滑样条或物理模型提取可辨识的趋势，再对趋势模型求导。平滑程度会引入偏差（你刻意忽略了一些高频变化），却抑制了不可恢复的高频噪声。偏差—方差的权衡在此以另一种形态出现：越平滑，噪声抑制越好，但被误归为噪声的真实局部变化也越多。

\subsection{complex-step：绕开实数相消的迂回策略}

如果被求导函数 \(f\) 可以解析延拓到复平面，并且你的全部代码路径都支持复数运算，有一个巧妙的迂回路线：

\[
f(x+ih)=f(x)+ihf'(x)-\frac{h^2}{2}f''(x)+O(h^3).
\]

取虚部并除以 \(h\)：

\[
f'(x)\approx\frac{\operatorname{Im}f(x+ih)}{h}.
\]

注意这个公式里没有两个接近实数相减的操作——\(f(x+ih)\) 和 \(f(x)\) 不在同一条实轴上竞争有效位。因此 \(h\) 可以取到 \(10^{-100}\) 甚至机器下溢的极限，截断误差被压到机器精度的平方级别，却不会被舍入相消的 U 形曲线干扰。

但代价同样明确。第一，函数必须可解析延拓——绝对值、条件分支、max/min、取整、逻辑判断、非解析特殊函数，任何一处不可解析都会让回答失去意义。第二，所有被 \(f\) 调用的子程序、外部库、乃至底层数学原语都必须支持复数运算并保持解析性。第三，复数路径的成本大约是实数路径的 2 到 8 倍——函数体越长，开销越大。这不是一个"对所有黑盒函数按下就行的通用技巧"。

另一种避开有限差分的路线是自动微分（automatic differentiation）。它通过对程序基本运算的链式法则传播，累积出精确的局部导数值——既没有截断误差，也没有步长选择的困扰。但它要求程序由可微基本操作构成，且仍然受原始程序本身的浮点舍入和不可微点（如分支条件跳变）影响。有限差分、复步微分和自动微分各自的适用边界，应作为工程选型时的一项清醒判断。
